% GENERATED FILE. Edit canonical lessons, then run npm run generate:books.
% canonical-source-hash: fnv1a64:2e4a1769b470c35c
% canonical-lessons: SA-C184-duratah, SA-C184-nikate, SA-C184-antara, SA-C184-uccaih, SA-C184-nicaih

\chapter{From afar, Close by, Between, Loudly, Softly}
\label{ch:sa-from-afar-close-by-between-loudly-softly}
\hlchaptermodality{184}

\begin{chapteropening}
\textbf{By the end of this chapter:} \emph{I can say from afar, close by, between, loudly and softly.}

Complete the last lesson of chapter 184: I can say from afar, close by, between, loudly and softly.
\end{chapteropening}

\section[\texorpdfstring{\sk{दूरतः} (dūrataḥ)}{दूरतः (dūrataḥ)}]{\sk{दूरतः} (dūrataḥ) — from afar}

\label{lesson:SA-C184-duratah}



\begin{quote}
Before the new one: say the Sanskrit for for a long time, then the Sanskrit for after a long time.
\end{quote}

\subsection*{You'll want to know: \sk{दूरतः}}
\textbf{\sk{दूरतः}} — \emph{dūrataḥ} — \textquotedblleft{}from afar\textquotedblright{}.

\begin{grammarlens}[title={What you've built}]
One more word for this chapter.
\end{grammarlens}

\subsection*{Your turn}
\begin{itemize}
\raggedright
  \item \emph{Say it:} \emph{dūrataḥ}
  \item \emph{Say it:} \emph{dūrataḥ}, once more
  \item \emph{Say it:} say \emph{dūrataḥ}
  \item \emph{Recall:} say the Sanskrit for for a long time, then the Sanskrit for after a long time, then say \emph{dūrataḥ} again
\end{itemize}

\begin{tcolorbox}[breakable,colback=teal!4,colframe=teal!35!black,title={Before you move on}]
What does \sk{दूरतः} mean? (\textquotedblleft{}From afar\textquotedblright{}.) Say it once more.
\end{tcolorbox}
\section[\texorpdfstring{\sk{निकटे} (nikaṭe)}{निकटे (nikaṭe)}]{\sk{निकटे} (nikaṭe) — close by, near}

\label{lesson:SA-C184-nikate}



\begin{quote}
Before the new one: say the Sanskrit for after a long time, then the Sanskrit for from afar.
\end{quote}

\subsection*{You'll want to know: \sk{निकटे}}
\textbf{\sk{निकटे}} — \emph{nikaṭe} — \textquotedblleft{}close by, near\textquotedblright{}.

\begin{grammarlens}[title={What you've built}]
One more word for this chapter.
\end{grammarlens}

\subsection*{Your turn}
\begin{itemize}
\raggedright
  \item \emph{Say it:} \emph{nikaṭe}
  \item \emph{Say it:} \emph{nikaṭe}, once more
  \item \emph{Say it:} say \emph{nikaṭe}
  \item \emph{Recall:} say the Sanskrit for after a long time, then the Sanskrit for from afar, then say \emph{nikaṭe} again
\end{itemize}

\begin{tcolorbox}[breakable,colback=teal!4,colframe=teal!35!black,title={Before you move on}]
What does \sk{निकटे} mean? (\textquotedblleft{}Close by, near\textquotedblright{}.) Say it once more.
\end{tcolorbox}
\section[\texorpdfstring{\sk{अन्तरा} (antarā)}{अन्तरा (antarā)}]{\sk{अन्तरा} (antarā) — between}

\label{lesson:SA-C184-antara}



\begin{quote}
Before the new one: say the Sanskrit for from afar, then the Sanskrit for close by.
\end{quote}

\subsection*{You'll want to know: \sk{अन्तरा}}
\textbf{\sk{अन्तरा}} — \emph{antarā} — \textquotedblleft{}between\textquotedblright{}.

\begin{grammarlens}[title={What you've built}]
One more word for this chapter.
\end{grammarlens}

\subsection*{Your turn}
\begin{itemize}
\raggedright
  \item \emph{Say it:} \emph{antarā}
  \item \emph{Say it:} \emph{antarā}, once more
  \item \emph{Say it:} say \emph{antarā}
  \item \emph{Recall:} say the Sanskrit for from afar, then the Sanskrit for close by, then say \emph{antarā} again
\end{itemize}

\begin{tcolorbox}[breakable,colback=teal!4,colframe=teal!35!black,title={Before you move on}]
What does \sk{अन्तरा} mean? (\textquotedblleft{}Between\textquotedblright{}.) Say it once more.
\end{tcolorbox}
\section[\texorpdfstring{\sk{उच्चैः} (uccaiḥ)}{उच्चैः (uccaiḥ)}]{\sk{उच्चैः} (uccaiḥ) — loudly, high}

\label{lesson:SA-C184-uccaih}



\begin{quote}
Before the new one: say the Sanskrit for close by, then the Sanskrit for between.
\end{quote}

\subsection*{You'll want to know: \sk{उच्चैः}}
\textbf{\sk{उच्चैः}} — \emph{uccaiḥ} — \textquotedblleft{}loudly, high\textquotedblright{}.

\begin{grammarlens}[title={What you've built}]
One more word for this chapter.
\end{grammarlens}

\subsection*{Your turn}
\begin{itemize}
\raggedright
  \item \emph{Say it:} \emph{uccaiḥ}
  \item \emph{Say it:} \emph{uccaiḥ}, once more
  \item \emph{Say it:} say \emph{uccaiḥ}
  \item \emph{Recall:} say the Sanskrit for close by, then the Sanskrit for between, then say \emph{uccaiḥ} again
\end{itemize}

\begin{tcolorbox}[breakable,colback=teal!4,colframe=teal!35!black,title={Before you move on}]
What does \sk{उच्चैः} mean? (\textquotedblleft{}Loudly, high\textquotedblright{}.) Say it once more.
\end{tcolorbox}
\section[\texorpdfstring{\sk{नीचैः} (nīcaiḥ)}{नीचैः (nīcaiḥ)}]{\sk{नीचैः} (nīcaiḥ) — softly, low}

\label{lesson:SA-C184-nicaih}



\begin{quote}
Before the new one: say the Sanskrit for between, then the Sanskrit for loudly.
\end{quote}

\subsection*{You'll want to know: \sk{नीचैः}}
\textbf{\sk{नीचैः}} — \emph{nīcaiḥ} — \textquotedblleft{}softly, low\textquotedblright{}.

\begin{grammarlens}[title={What you've built}]
One more word for this chapter.
\end{grammarlens}

\subsection*{Your turn}
\begin{itemize}
\raggedright
  \item \emph{Say it:} \emph{nīcaiḥ}
  \item \emph{Say it:} \emph{nīcaiḥ}, once more
  \item \emph{Say it:} say \emph{nīcaiḥ}
  \item \emph{Recall:} say the Sanskrit for between, then the Sanskrit for loudly, then say \emph{nīcaiḥ} again
\end{itemize}

\begin{tcolorbox}[breakable,colback=teal!4,colframe=teal!35!black,title={Before you move on}]
What does \sk{नीचैः} mean? (\textquotedblleft{}Softly, low\textquotedblright{}.) Say it once more.
\end{tcolorbox}
